\chapter{Introduction}
Il mondo attuale presenta molte sfide relative al cambiamento climatico, nello specifico questo studio si preoccuperà del sempre maggiore dispendio di risorse idriche da utilizzare per le attività umane ed al contempo la sempre maggiore scarsità di risorse, soprattutto in zone del mondo particolarmente esposte a fenomeni di desertificazione o di sistemi idrici non efficienti.\\
Fra i vari campi umani più dispendiosi a livello di consumo d'acqua l'agricoltura è sicuramente uno di questi arrivando a consumare cifre tra il 70\% e l'80\% delle risorse d'acqua dolce utilizzata con picchi nei paesi in via di sviluppo. Un parametro, che è stato evidenziato come di fondamentale rilevanza in questo settore da Elbeltagi et al. ~\cite{elbeltagi2023vpd} , è il  deficit di pressione di vapore (VPD), valore che rappresenta la "forza" con cui l'aria "attira" l'umidità fuori dalle piante,un aumento di tale misura porta  le piante ad una situazione di siccità atmosferica anche in presenza di terreno perfettamente bagnato, come ben dimostrato da Novick et al 2024 .~\cite{novick2024vpd}, il VPD "ha effetti negativi e a cascata su quasi tutti gli aspetti della funzionalità delle piante, tra cui la fotosintesi, lo stato idrico, la crescita e la sopravvivenza", inoltre esso influenza anche l'assorbimento di carbonio, la biodiversità, la disponibilità di risorse idriche e le rese agricole; Questo studio si concentrerà proprio su di esso e nello specifico sulla correlazione che il VPD ha con questi due ultimi fattori.\\
l'aumento del valore del VPD è un fenomeno globale, ma nonostante ciò la ricerca su come integrare questi impatti in regimi di gestione resilienti è ancora in una fase iniziale (Novick et al 2024 .~\cite{novick2024vpd}), questa tesi in linea con quanto citato vuole collocarsi nel novero delle possibili ricerche che si pongono come obiettivo quello di trovare strumenti di gestione resilienti proprio atti all'adattabilità in uno scenario di VPD in costante crescita negli anni, ovvero lo scenario che ci ritroviamo ad affrontare. Un aumento generalizzato del VPD può quindi portare ad un maggiore utilizzo dell'acqua a causa di, per esempio, tecniche di nebulizzazione per rendere l'aria meno secca e quindi meno aggressiva per quanto riguarda il consumo della stessa dalle piante, di conseguenza l'aumento di questo parametro può, ancor più che una variazione delle precipitazioni, fenomeno che comunque se verificatosi in concomitanza con un aumento del VPD può esacerbare i danni dello stesso (Will et al.~\cite{will2013vulnerability} , 2013 ; Duan et al.~\cite{duan2014elevated}, 2014 ), avere un impatto significativo sulla resa e produttività delle piante e, di conseguenza, dei raccolti. Essendo il VPD un parametro fondamentale oltre che per il prodotto delle piante anche per la crescita delle stesse, in quanto è stato dimostrato che le piante crescono maggiormente durante la notte ovvero quando i valori di VPD sono più bassi(Zweifel et al.~\cite{zweifel2021night},  2021), perciò un aumento dello stesso potrebbe portare, come sta già facendo, ad un incremento della mortalità forestale legata alla siccità (Yuan et al.~\cite{yuan2019global}., 2019, fenomeno ulteriormente aggravato dall'inaridimento del suolo, evento fortemente influenzato ancora una volta dal crescente VPD (Li et al.~\cite{li2023soil}  2023). \\
Da ciò si può vedere come un aumento generalizzato del VPD potrà portare ad una maggiore difficoltà nel trattenere l'acqua nelle piante ed ad una maggiore mortalità delle stesse causata da fenomeni di disidratazione e carenza di carbonio, portando in generale alla Terra nella sua interezza ad attraversare una maggiore siccità atmosferica a livello globale, fenomeno che può solo andare peggiorando dalle condizioni climatiche in progressivo peggioramento da decenni, per esempio dall'aumento della temperatura superficiale globale di circa 0,2 °C per decennio (IPCC,~\cite{ipcc2019report} 2019 ), tenendo presente ciò l'equazione di Clausius-Clapeyron (Bolton,~\cite{bolton1980computation} 1980; Iribarne e Godson, ~\cite{iribarne1981atmospheric} 1981) ci dice che  la pressione del vapore acqueo saturo aumenterà di circa il 7\% per ogni grado kelvin di aumento della temperatura atmosferica e per far sì che questo non provochi squilibri all'ambiente come per esempio un aumento del VPD, ciò dovrebbe accadere in concomitanza di un aumento della concentrazione di vapore acqueo atmosferico, ovvero dell'umidità dell'aria; Lo studio di Simpson et al.~\cite{simpson2023arid} (2023) però mostra come per l'appunto, soprattutto nelle zone aride e semi-aride, anche se ciò vale in modo generalizzato in tutto il mondo nelle stagioni calde, la concentrazione di vapore acqueo atmosferico non stia aumentando ,causando tra le altre cose, un aumento del VPD (a testimonianza di ciò si può vedere già adesso come nelle zone in cui già tutt'ora si sia verificato un aumento del vpd e della temperatura, che come si è visto hanno una forte relazione tra di loro, esso sia stato un fatto molto limitante per la crescita e la produttività degli alberi nelle piantagioni, Williams et al.~\cite{williams2013forest} . , 2013 ; Babst et al.~\cite{babst2019twentieth} , 2019).\\
I fenomeni correlati al vpd hanno un effetto ancora maggiore se si osservano contesti di coltivazione in serra , che non riguardano una parte risibile del totale delle piantagioni, occupando a livello mondiale circa 400000 ettari di terreno(Baudoin,~\cite{baudoin2020greenhouses} 2020), in questa tipologia di colture la gestione del VPD all'interno della serra è essenziale per mantenere alta la produttività delle piante, si è però riscontrato anche come il risparmio idrico effettivo si abbia solo in colture con una densità di impianto abbastanza estesa, in quanto nelle colture con una densità troppo bassa il risparmio d'acqua verificatosi nella minore quantità della stessa utilizzata nell'effettiva irrigazione viene compensato da quella utilizzata per la nebulizzazione, o altre tecniche atte alla modifica del VPD all'interno delle serre (Zhang et al.~\cite{zhang2017vpd} 2017), ciò può suggerire anche come sia essenziale per un efficientamento ed una salvaguardia delle risorse ammodernare ed ottimizzare le colture esistenti in modo da poter sfruttare i vantaggi dimostrati dallo studio e poter anche applicare ottimamente i risultati di questa tesi; In conclusione per un sempre maggiore efficientamento nell'utilizzo di risorse e nella produttività il VPD si dimostra ancora un parametro fondamentale da tenere in considerazione, permettendo un'efficace moderazione dello stress idrico delle piante, un aumento della produttività delle stesse ed infine un'ottimizzazione delle risorse utilizzate in colture con una densità d'impianto abbastanza alta.\\
Oltre che mostrare come il vpd fosse un parametro fondamentale per una gestione efficiente delle coltivazioni Elbeltagi et al.~\cite{elbeltagi2023vpd}hanno evidenziato anche come esso sia ancora più fondamentale in paesi a rischio desertificazione e/o con grandi problemi nella disponibilità di acqua per l'irrigazione, problemi che andranno sempre più esacerbandosi col riscaldamento climatico in atto, nello specifico quelli del Nord Africa e della penisola arabica, indi per cui riuscire a prevedere i futuri valori di VPD può essere fondamentale per programmare in anticipo la quantità d'acqua di cui una determinata piantagione avrà bisogno per evitare situazioni di sovra-irrigazione o di sotto-irrigazione, mantenendo  quindi la massima produttività col minimo dispendio di risorse.\\
Nello studio Elbeltagi et al.~\cite{elbeltagi2023vpd} hanno utilizzato sei diversi algoritmi di machine learning implementati attraverso il software WEKA, i sei modelli erano i seguenti:  Linear Regression (LR), additive regression trees (BART), Random SubSpace (RSS), Random Forest (RF), REPTree, e M5P, i quali sono stati utilizzati per realizzare previsioni sul vpd a lungo termine su sei regioni della foce del Nilo, i diversi risultati dei 6 modelli sono poi stati paragonati al fine di verificare se l'approccio attraverso modelli di machine learning per questo tipo di previsione fosse quello corretto, ed in caso affermativo quale di essi fosse il modello dalle performance migliori. E' stato riscontrato come l'approccio predittivo basato su modelli di machine learning abbia dato esito positivo con dei tassi di accuratezza abbastanza elevati da essere ritenuti affidabili, da ciò rilevando come i problemi di desertificazione e di carenza idrica, a cui sono sottoposti i paesi presi in esame nello studio, possono essere molto simili a quelli che la Sicilia sta iniziando ad affrontare, come viene attestato dallo studio World Weather Attribution 2024.~\cite{wwa2024sicily}, la quale ricerca attesta come gli eventi straordinari quali per esempio le straordinarie ondate di calore che abbiamo osservato negli ultimi anni siano state molto più forti di quanto sarebbero dovute essere senza gli effetti del cambiamento climatico, inoltre attesta anche come "i valori osservati per l'evapotraspirazione potenziale e la temperatura, sarebbero stati pressoché impossibili da raggiungere senza i cambiamenti climatici di origine antropica. Si conclude pertanto che questo aumento della gravità della siccità è dovuto principalmente al forte incremento delle temperature estreme causato dai cambiamenti climatici di origine antropica.", ad aggiungersi a questo abbiamo tutti sotto gli occhi i problemi che la rete idrica siciliana ha nel soddisfare i bisogni dei suoi abitanti e lavoratori, da ciò si evince come una migliore gestione delle risorse idriche sia essenziale non solo per il settore agricolo siciliano ma in particolar modo per esso.\\
Come la rilevanza del VPD e quindi della sua previsione sia uno strumento fondamentale per l'efficientamento dell'irrigazione e quindi di particolare utilità nel contesto siciliano è dimostrato dalla strettissima relazione presente tra esso e il potenziale idrico fogliare (lwp, leaf water potential), misura che indica lo stato idrico di una pianta, essenziale per capire quanta energia serva alla pianta per trattenere o muovere l'acqua al suo interno, questa relazione è ancora più forte di quella tra il lwp e la temperatura, rendendo così il vpd uno strumento essenziale nell'organizzazione dell'irrigazione, con l'obiettivo di utilizzare la giusta quantità di acqua al momento giusto , richiesta dal raccolto, così da poter minimizzare le risorse usate e massimizzare la produzione.
L'obiettivo di questo studio è quindi quello di applicare l'ipotesi alla base dello studio di Elbeltagi et al.~\cite{elbeltagi2023vpd} al caso siciliano, così da verificare se i modelli di machine learning e deep learning testati, ovvero  Linear Regression (LR), Additive Regression trees (AR), Random SubSpace (RSS), Random Forest (RF), REPTree, M5P, Long Short-Term Memory (LSTM) e Gated recurrent unit (GRU), riescano a captare le specificità delle varie province siciliane, ed inoltre capire attraverso comparazioni tra gli stessi quale di essi sia la scelta corretta a tale scopo, permettendo una corretta previsione del vpd per consentire agli agricoltori siciliani una più facile ed automatizzata previsione del fabbisogno idrico delle loro colture evitando una sovra-irrigazione nei casi di vpd basso o una sotto-irrigazione in casi di vpd sopra la media.
\label{sec:introduction}